\makeatletter
\def\tkz@KiviatGrad[#1](#2){% 
\begingroup
\pgfkeys{/kiviatgrad/.cd,
graduation distance= 0 pt,
prefix ={},
suffix={},
unity=1,
% added the following four lines
label precision/.store in=\gradlabel@precision,
label precision=1,
zero point/.store in=\tkz@grad@zero,
zero point=0
}
\pgfqkeys{/kiviatgrad}{#1}% 
\let\tikz@label@distance@tmp\tikz@label@distance
\global\let\tikz@label@distance\tkz@kiv@grad
  % started loop at zero, to get label for center point
 \foreach \nv in {0,...,\tkz@kiv@lattice}{
 % skipped the \pgfmathparse, moved
 % \tkz@kiv@unity*\nv-\tkz@grad@zero
 % into \pgfmathsetmacro. The last term is added
 \pgfmathsetmacro{\result}{\tkz@kiv@unity*\nv+20} % change from \pgfmathtruncatemacro
 % NOTE: use \gradlabel@precision for the precision
 \protected@edef\tkz@kiv@gd{%
    \tkz@kiv@prefix%
    \pgfmathprintnumber[precision=\gradlabel@precision,fixed]{\result}%% used \pgfmathprintnumber instead of "$\result$"
    \tkz@kiv@suffix} 
    \path[/kiviatgrad/.cd,#1] (0:0)--(360/\tkz@kiv@radial*#2:\nv*\tkz@kiv@gap)
       % NOTE: (360/\tkz@kiv@radial*#2)-90  -> (360/\tkz@kiv@radial*#2)
       % removed 90 because you rotate all the diagrams
       node[label=(360/\tkz@kiv@radial*#2):\tiny\tkz@kiv@gd] {}; % added \tiny
      }
 \let\tikz@label@distance\tikz@label@distance@tmp  
\endgroup
}%


\def\tkz@KiviatLine[#1](#2,#3){% 
\begingroup
\pgfkeys{/kiviatline/.cd,
fill= {},
opacity=.5,
% add these two lines, similar to for kiviatline
zero point/.store in=\tkz@line@zero,
zero point=0
}
%
% the code below has a lot of polar coordinates, of the form
%   (angle: <value> * <some factor>)
% all have been modified to get them on the form
%   (angle:{(<value>+\tkz@line@zero) * <some factor>})
% note extra braces, which are required when a component contains ()
% 
\pgfqkeys{/kiviatline}{#1}%   opacity ??????
\ifx\tkzutil@empty\tkz@kivl@fill \else 
\begin{scope}[on background layer]
 \path[fill=\tkz@kivl@fill,opacity=\tkz@kivl@opacity] (360/\tkz@kiv@radial*0:{(#2+\tkz@line@zero)*\tkz@kiv@gap*\tkz@kiv@step})   
\foreach \v [count=\rang from 1] in {#3}{%  
 -- (360/\tkz@kiv@radial*\rang:{(\v+\tkz@line@zero)*\tkz@kiv@gap*\tkz@kiv@step}) } -- (360/\tkz@kiv@radial*0:{(#2+\tkz@line@zero)*\tkz@kiv@gap*\tkz@kiv@step}); 
 \end{scope}
 \fi       
% added overlay option because something weird happened with the bounding box
\draw[#1,opacity=1,overlay] (0:{(#2+\tkz@line@zero)*\tkz@kiv@gap}) plot coordinates {(360/\tkz@kiv@radial*0:{(#2+\tkz@line@zero)*\tkz@kiv@gap*\tkz@kiv@step})}  
\foreach \v [count=\rang from 1] in {#3}{%  
 -- (360/\tkz@kiv@radial*\rang:{(\v+\tkz@line@zero)*\tkz@kiv@gap*\tkz@kiv@step}) plot coordinates {(360/\tkz@kiv@radial*\rang:{(\v+\tkz@line@zero)*\tkz@kiv@gap*\tkz@kiv@step})}} -- (360/\tkz@kiv@radial*0:{(#2+\tkz@line@zero)*\tkz@kiv@gap*\tkz@kiv@step});   
\endgroup
}%  

\makeatother
%%%%%%%%%%%%%%%%%%%%%%%%%%%%%%%
\definecolor{c1}{HTML}{81BECE}
\definecolor{c2}{HTML}{378BA4}
%%%%%%%%%%%%%%%%%%%%%%%%%%%%%%%
\tikzset{global scale/.style={
    scale=#1,
    every node/.append style={scale=#1}
  }
}
%%%%%%%%%%%%%%%%%%%%%%%%%%%%%%%%%%%%%%%%%%
\begin{tikzpicture}[
  label distance=13em,
  global scale = 0.73,
  plot1/.style={
    thin,
    draw=c1,
    fill=c1,
    mark=square*,
    mark options={
     ball color=blue, % note ball color inside mark options
     mark size=4pt
    },
    opacity=.5
  },
  plot2/.style={
    thin,
    draw=c2,
    fill=c2,
    mark=triangle*,
    mark options={
      mark size=4pt,
      ball color=blue
    },
    opacity=.5
  },
  plot3/.style={
    thin,
    draw=red!30,
    fill=red!30,
    mark=square*,
    mark options={
      mark size=4pt,
      ball color=yellow
    },
    opacity=.5
  },
  plot4/.style={
    thin,
    draw=red!50,
    fill=red!50,
    mark=triangle*,
    mark options={
      mark size=4pt,
      ball color=blue
    },
    opacity=.5
  }
]
%%%%%%%%%%%%%%%%设置画布大小%%%%%%%%%%%%%%%%%%%%%%%%
\useasboundingbox (-25em,-5em) rectangle (10em,7em);
% Correlation coefficient
\begin{scope}[scale=0.36,local bounding box=a,shift={(-50em,0)}]

% define some macros for convenience
\newcommand\KivStep{0.1}
\pgfmathsetmacro\Unity{1/\KivStep}
\newcommand\zeroshift{20.00}

 \tkzKiviatDiagram[
   radial  style/.style ={-},
   rotate=90, 
   lattice style/.style ={black!30},
   step=\KivStep,
   gap=1,
   lattice=5,
]% 
% de$\rightarrow$en  en$\rightarrow$de
{cs$\rightarrow$en,cs$\rightarrow$it,,ro$\rightarrow$en,en$\rightarrow$cs,it$\rightarrow$cs,,en$\rightarrow$ro}
\node [align=center,rotate=90]at (0em,16em){de$\rightarrow$en};
\node [align=center,rotate=270]at (0em,-16em){en$\rightarrow$de};

% with tm
\tkzKiviatLine[
  % /kiviatline/zero point=\zeroshift,
  plot3
](38.58,27.93,39.74,39.66,32.58,27.22,33.56,30.18)
% wo tm
\tkzKiviatLine[
  % /kiviatline/zero point=\zeroshift,
  plot4
](17.50,7.02,21.37,16.34,8.02,2.85,7.06,8.74)

\tkzKiviatGrad[unity=\Unity, label precision=2, zero point=\zeroshift](0) % set unity as 0.1
\end{scope}
%PBIAS
\begin{scope}[scale=0.36,local bounding box=b,shift={(-15em,0)}]

\newcommand\KivStep{0.1}
\pgfmathsetmacro\Unity{1/\KivStep}
\newcommand\zeroshift{20.00}

 \tkzKiviatDiagram[
   radial  style/.style ={-},
   rotate=90,
   lattice style/.style ={black!30},
   step=\KivStep,
   gap=1,
   lattice=5
]% 
% it$\rightarrow$en en$\rightarrow$it
{es$\rightarrow$en,fr$\rightarrow$en,,de$\rightarrow$fr,en$\rightarrow$es,en$\rightarrow$fr,,fr$\rightarrow$de}
\node [align=center,rotate=90]at (0em,16.5em){it$\rightarrow$en};
\node [align=center,rotate=270]at (0em,-16.5em){en$\rightarrow$it};
% with tm
\tkzKiviatLine[
  % /kiviatline/zero point=\zeroshift,
plot1
](41.25,44.45,41.80,33.52,39.18,44.60,37.71,30.08)
% wo tm
\tkzKiviatLine[
  % /kiviatline/zero point=\zeroshift,
plot2
](21.89,24.06,21.13,14.78,17.10,23.78,13.53,8.93)
%%%%%%%%
\tkzKiviatGrad[unity=\Unity, zero point=\zeroshift](0) % set unity as 0.1
\end{scope}

\end{tikzpicture}